\documentclass[12pt]{article}
\usepackage[utf8]{inputenc}
\usepackage[T1]{fontenc}
\usepackage[spanish]{babel}
\usepackage{geometry}
\usepackage{amsmath,amssymb}
\usepackage{enumitem}

\geometry{margin=2.5cm}

\begin{document}

\begin{center}
{\Large Transcripción cruda}\\[2mm]
{\large 25May26}
\end{center}

\section*{Fotos 1--2}

\textit{Teo.} Sea $(X,d_1)$ un espacio métrico. Sea
\[
d_2:X\times X\to[0,\infty),\qquad d_2(x,y)=\frac{d_1(x,y)}{1+d_1(x,y)}.
\]
Entonces:
\begin{enumerate}[label=\alph*)]
    \item $d_2$ es métrica.
    \item $d_1$ y $d_2$ son equivalentes, es decir, tienen los mismos conjuntos abiertos.
\end{enumerate}

\textit{Dem.} a) Ya lo sabemos. b) Solo probaremos b).
Sea
\[
\mathcal F_i=\{A\subset X\mid A \text{ es abierto en }(X,d_i)\},\qquad i=1,2.
\]
P.D. $\mathcal F_1=\mathcal F_2$.

Primero probaremos que $\mathcal F_2\subset \mathcal F_1$.
Usaremos que $d_2(x,y)\le d_1(x,y)$.
Sea $A\in\mathcal F_2$. Sea $a\in A$, $\exists r>0$ t.q.
\[
N_r^{d_2}(a)\subset A.
\]
Si $b\in N_r^{d_1}(a)$, entonces
\[
d_2(a,b)\le d_1(a,b)<r,
\]
luego $b\in N_r^{d_2}(a)$.
Así,
\[
N_r^{d_1}(a)\subset N_r^{d_2}(a)\subset A,
\]
de modo que $A\in\mathcal F_1$.

Recíprocamente, sea $A\in\mathcal F_1$.
Sea $a\in A$, $\exists r>0$ t.q.
\[
N_r^{d_1}(a)\subset A.
\]
Sea $f:[0,1)\to[0,\infty)$, $f(z)=\dfrac{z}{1-z}$.
Observemos que
\[
f(d_2(x,y))=\frac{d_2(x,y)}{1-d_2(x,y)}=d_1(x,y),
\]
porque $0\le d_2(x,y)<1$.
Además, $f$ es continua en $0$.
Entonces existe $\delta_r>0$ tal que
\[
0\le z<\delta_r \implies 0\le f(z)<r.
\]
Si $b\in N_{\delta_r}^{d_2}(a)$, entonces
\[
0\le d_2(a,b)<\delta_r
\]
y por tanto
\[
0\le f(d_2(a,b))=d_1(a,b)<r,
\]
así que $b\in N_r^{d_1}(a)\subset A$.
Luego
\[
N_{\delta_r}^{d_2}(a)\subset A,
\]
y $A\in\mathcal F_2$.
Concluimos que $\mathcal F_1=\mathcal F_2$.

\section*{Fotos 3--10}

\textit{Teo.} Sean $a,b,c,d\in\mathbb R$ tales que $a<b$ y $c<d$.
Sea
\[
J_0=[a,b]\times[c,d]\subset\mathbb R^2.
\]
Entonces $J_0$ es compacto.

\textit{Dem.} Supóngase que $J_0$ no es compacto.
Por tanto, existe una cubierta abierta $\mathcal U=\{U_\alpha\}_{\alpha\in\Lambda}$ de $J_0$ de la cual no se puede extraer una subcubierta finita.

Dividimos $J_0$ en cuatro rectángulos cerrados de área igual a $\frac14$ del área de $J_0$:
\[
[a,m_0]\times[c,n_0],\quad [m_0,b]\times[c,n_0],\quad [a,m_0]\times[n_0,d],\quad [m_0,b]\times[n_0,d],
\]
donde
\[
m_0=\frac{a+b}{2},\qquad n_0=\frac{c+d}{2}.
\]
Al menos uno de ellos no puede cubrirse con un número finito de elementos de $\mathcal U$; lo denotamos por
\[
J_1=[m_1,n_1]\times[\sigma_1,\rho_1]\subset J_0.
\]
Repetimos el procedimiento y obtenemos una sucesión decreciente de rectángulos cerrados
\[
J_0\supset J_1\supset J_2\supset\cdots
\]
tal que cada $J_k$ no puede cubrirse con una subfamilia finita de $\mathcal U$.

Escribiendo
\[
J_k=[m_k,n_k]\times[\sigma_k,\rho_k],
\]
se tiene
\[
m_k\le m_{k+1}\le x\le n_{k+1}\le n_k,\qquad
\sigma_k\le \sigma_{k+1}\le y\le \rho_{k+1}\le \rho_k.
\]
En particular,
\[
\lim_{k\to\infty}(n_k-m_k)=0,\qquad \lim_{k\to\infty}(\rho_k-\sigma_k)=0.
\]
Luego existe un único punto
\[
(x,y)\in\bigcap_{k=0}^\infty J_k.
\]

Como $\mathcal U$ cubre a $J_0$, existe $U_{\alpha_0}\in\mathcal U$ tal que $(x,y)\in U_{\alpha_0}$.
Puesto que $U_{\alpha_0}$ es abierto, existe $\varepsilon>0$ tal que
\[
N_\varepsilon^{d_2}((x,y))\subset U_{\alpha_0},
\]
donde
\[
d_2((x,y),(u,v))=\max\{|x-u|,|y-v|\}.
\]
Entonces, para $k$ grande, $J_k\subset N_\varepsilon^{d_2}((x,y))\subset U_{\alpha_0}$, contradicción.
Por tanto, $J_0$ es compacto.

\section*{Fotos 11--21}

\textit{Teo.} Sea $K\subset\mathbb R^d$. Entonces son equivalentes:
\begin{enumerate}[label=\alph*)]
    \item $K$ es cerrado y acotado.
    \item $K$ es compacto.
    \item Todo subconjunto infinito de $K$ tiene punto límite en $K$.
\end{enumerate}

\textit{Dem.} Esquema: $a)\Rightarrow b)\Rightarrow c)\Rightarrow a)$.

\textit{$a)\Rightarrow b)$.} Supongamos que $K$ es cerrado y acotado.
Entonces $K$ cabe en una celda-$d$,
\[
K\subset [a_1,b_1]\times\cdots\times[a_d,b_d],
\]
con $a_i<b_i$.
Ya sabemos que dicha celda es compacta.
Como $K$ es cerrado en un compacto, resulta que $K$ es compacto.

\textit{$b)\Rightarrow c)$.} Sea $L\subset K$ infinito.
Por compacidad de $K$, todo subconjunto infinito de $K$ tiene punto límite en $K$.
Luego $L$ tiene punto límite en $K$.

\textit{$c)\Rightarrow a)$.} Falta probar que $K$ es cerrado y acotado.
Si $K$ no es acotado, entonces para cada $n\in\mathbb N$ existe $x_n\in K$ tal que
\[
\|x_n\|=d(x_n,0)>n.
\]
Sea
\[
S=\{x_1,x_2,\dots\}\subset K.
\]
El conjunto $S$ es infinito.
Si $y\in S'$, entonces para todo $N\in\mathbb N$ existe $z\in S\setminus\{y\}$ tal que
\[
d(z,y)<\frac1N.
\]
Eso fuerza que haya elementos de $S$ arbitrariamente cercanos a $y$.
Pero como $\|x_n\|>n$, la sucesión no puede acumularse en un punto de $\mathbb R^d$.
Por tanto $K$ debe ser acotado.

Si $K$ no es cerrado, existe $x_0\in \overline K\setminus K$.
Entonces hay una sucesión $(x_n)\subset K$ con $x_n\to x_0$.
El conjunto $S=\{x_1,x_2,\dots\}$ es infinito y su único punto límite es $x_0$.
Como $x_0\notin K$, esto contradice $c)$.
Luego $K$ es cerrado.

Concluimos que $K$ es cerrado y acotado.

\end{document}
